\chapter*{Preface}
\addcontentsline{toc}{chapter}{Preface}
\markboth{PREFACE}{PREFACE}

This is a course about one problem:
\[
  \min_{x \in \R^n} \; \tfrac12 x\T G x - a\T x
  \qquad\text{subject to}\qquad C\T x \ge b,
\]
a convex quadratic minimised over a polyhedron. That is a narrow subject to
give a book to, and the narrowness is the point. The problem is small enough
that everything about it can be proved, and large enough that the proving
teaches most of what a graduate student needs to know about constrained
optimization: optimality conditions and when they are sufficient, duality,
certificates, active sets and their combinatorics, saddle-point linear algebra,
Krylov subspace methods, conditioning, degeneracy, warm starts, and how to tell
whether a method is fast for a reason or by accident.

It is also a problem the reader has almost certainly met already, under some
other name. Mean--variance portfolio selection is exactly this problem. So is
non-negative least squares. So is every step of a sequential quadratic
programming method, every horizon of a linear model predictive controller, the
dual of a support vector machine, and---after one substitution that
Chapter~\ref{ch:homotopy} makes precise---the LASSO. Four literatures developed
solvers for it in parallel for half a century, largely without reference to one
another. One of this book's aims is to make that duplication visible, and then
useful.

\section*{The organising question}

Nearly all the difficulty in the problem above is combinatorial rather than
analytic. If someone tells you \emph{which} constraints hold with equality at
the solution---the \emph{active set}---the rest is a linear system. Everything
else is a strategy for finding that set. The book is organised around four such
strategies, and they turn out to exhaust the interesting possibilities:

\begin{enumerate}
\item \textbf{Walk to it}, one constraint at a time, maintaining optimality
      conditions on one side and driving violation on the other. This is the
      dual active-set method of Goldfarb and Idnani (Chapter~\ref{ch:walking}).
\item \textbf{Guess it}, wholesale, and repair the guess from the signs that
      come back. This is the primal--dual active set strategy, equivalently
      block principal pivoting on a linear complementarity problem
      (Chapter~\ref{ch:guessing}).
\item \textbf{Trace it}, as a parameter moves, recording the finitely many
      values at which it changes. This is parametric quadratic programming, and
      it is Markowitz's Critical Line Algorithm, least angle regression, the
      support-vector-machine path, and explicit model predictive control, all at
      once (Chapters~\ref{ch:parametric}--\ref{ch:reading}).
\item \textbf{Refuse to form the matrix}, and reach the problem only through the
      three linear actions an active-set method actually performs
      (Chapters~\ref{ch:krylov} and~\ref{ch:operator}).
\end{enumerate}

Two of these strategies come with finite-termination theorems and two do not.
The book's central methodological claim, argued from Chapter~\ref{ch:certificate}
onward, is that this matters less than it appears, because strict convexity makes
the Karush--Kuhn--Tucker conditions \emph{sufficient}. A candidate solution can
therefore be \emph{certified} rather than trusted. A method that may fail, but
that can recognise its own success, needs a certificate and not a convergence
theory. Much of modern computational practice runs on exactly that trade, and it
is worth teaching deliberately rather than leaving students to infer it.

\section*{Who this is for}

The reader is assumed to have had a first course in linear algebra including the
QR and Cholesky factorisations, a course in analysis sufficient for convexity and
compactness arguments, and enough numerical analysis to know that floating-point
arithmetic is not real arithmetic. No prior optimization course is assumed;
Chapter~\ref{ch:certificate} develops the Karush--Kuhn--Tucker conditions from
scratch for this problem class, which is the one setting where they can be proved
in a page and are sufficient as well as necessary. No finance, statistics, or
control background is assumed either: the applications are introduced where they
are used, and Appendix~\ref{app:notation} gives a translation table between the
four dialects.

Programming is not required to read the book, but the material is written by
people who implement, and it shows. Where a formula is a specification rather than
an implementation, the text says so. Where a constant in an algorithm exists to
defeat a floating-point hazard rather than a mathematical one, the text says that
too. Chapters~\ref{ch:numerics} and~\ref{ch:measurement} are given over entirely
to that layer, on the view that a graduate course which stops at the theorem
leaves students to rediscover the hard half in their first job.

\section*{How to use it}

The five parts are cumulative, but not equally so.

\begin{itemize}
\item \textbf{Part I} (Chapters~\ref{ch:oneproblem}--\ref{ch:freeblock}) is the
      foundation and should be read in order. It sets the problem, proves the
      optimality conditions and the sufficiency that everything later leans on,
      develops the combinatorial geometry of active sets, and assembles the
      linear algebra of the free block.
\item \textbf{Part II} (Chapters~\ref{ch:walking}--\ref{ch:krylov}) gives the
      three single-point strategies. The three chapters are largely independent
      of one another and can be taken in any order.
\item \textbf{Part III} (Chapters~\ref{ch:parametric}--\ref{ch:reading}) is the
      parametric theory and the cross-field synthesis. It depends on Part~I but
      only lightly on Part~II.
\item \textbf{Part IV} (Chapters~\ref{ch:operator}--\ref{ch:warmstart}) is about
      structure: reading the quadratic form as an operator, conditioning the
      problem, and exploiting families of related problems. It draws on both
      Part~II and Part~III.
\item \textbf{Part V} (Chapters~\ref{ch:numerics}--\ref{ch:measurement}) is the
      practice: numerics, degeneracy, and honest measurement.
\end{itemize}

A one-semester course can be built from Parts~I and~II with
Chapters~\ref{ch:parametric}, \ref{ch:homotopy}, and~\ref{ch:operator} added, which
is roughly ten weeks of material. A course emphasising statistics or machine
learning should keep Part~III entire and may compress Chapter~\ref{ch:walking}. A
course emphasising numerical linear algebra should keep
Chapters~\ref{ch:freeblock}, \ref{ch:krylov}, \ref{ch:operator},
and~\ref{ch:conditioning}, and may treat Part~III as reading.

Exercises follow every chapter. They are of three kinds, not marked as such but
easy to tell apart: routine verifications that a claim in the text was
understood; derivations that extend a result to a case the text skipped; and
open-ended computational problems, which say ``implement'' and mean it. The last
kind are the ones worth setting.

\section*{Sources}

The book is assembled from a set of companion research papers, and a reader who
wants the primary sources, the full experimental evidence, or the exact
statements as their authors made them should go to them.
Chapters~\ref{ch:walking} and~\ref{ch:guessing} follow \emph{Goldfarb--Idnani
Revisited: Invariants, Certificates, and the Limits of Guessing}.
Chapter~\ref{ch:krylov} follows \emph{Non-Negative Conjugate Gradients}
\cite{schmelzer2026nncg}. Chapters~\ref{ch:operator} and~\ref{ch:conditioning}
draw on \emph{Matrix-Free Methods for Long-Only Portfolio Optimization}
\cite{schmelzer2026matrixfree} and \emph{From Marchenko--Pastur to Woodbury}
\cite{schmelzer2026rmt}. Part~III follows \emph{One Homotopy Across Statistics,
Optimization, and Finance} \cite{schmelzer2026perspective} and the \texttt{cvxcla}
software paper \cite{cvxcla}. Where the book states an empirical number, it is
one of theirs, and the citation says which.

The book's own contribution is the arrangement: one notation, one set of proofs,
and the argument that the four strategies of the organising question are answers
to a single question rather than four separate subjects. Any error in the
translation is the book's and not the papers'.

\bigskip
\noindent\emph{Abu Dhabi}
